\section{Introduction}
\label{sec:intro}


Fibonacci numbers $F_n$ and Lucas numbers $L_n$ are among the most widely studied objects in mathematics. They are defined recursively by
\begin{align*}
F_0 = 0, &\quad F_1 = 1,\quad F_{n} = F_{n-1} + F_{n-2} \quad(n \ge 2), \\
L_0 = 2, &\quad L_1 = 1,\quad L_{n} = L_{n-1} + L_{n-2} \quad(n \ge 2).
\end{align*}
Many authors have studied various arithmetic properties of these sequences.
For example, Brillhart--Montgomery--Silverman \cite{brillhart1988tables} studied factorizations of Fibonacci and Lucas numbers for small $n$. Dubner--Keller \cite{dubner1999new} studied Fibonacci and Lucas numbers that are prime, and it remains open whether there are infinitely many such numbers.
Some works are also interested in determining all perfect powers in the Fibonacci and Lucas sequences.
For instance, Cohn \cite{cohn1964square} gave an elementary proof showing that $F_0 = 0$, $F_1 = F_2 = 1$, and $F_{12} = 144$ are the only squares.
Bugeaud, Mignotte, and Siksek \cite{bugeaud2006classical} showed that $0$, $1$, $8$, and $144$ are the only perfect powers appearing in the Fibonacci sequence (see Section \ref{subsec:prelim_diophantine}).

An analogue of this problem can be formulated for \emph{Fibonacci polynomials} $F_n(T) \in \bZ[T]$, the polynomial analogues of Fibonacci numbers, defined by
\begin{equation}
\label{eqn:fibo_poly}
    F_0(T) = 0, \quad F_1(T) = 1, \quad F_{n}(T) = T F_{n-1}(T) + F_{n-2}(T).
\end{equation}
For example, several authors have studied the irreducibility of Fibonacci polynomials. Webb and Parberry showed that $F_n(T)$ is irreducible over $\bQ$ if and only if $n$ is prime \cite{webb1969divisibility}.
Bergum and Hoggatt considered the analogous problem for Lucas polynomials and their generalizations in \cite{bergum1974irreducibility}.
In \cite{kitayama2017irreducibility}, Kitayama and Shiomi studied irreducibility over finite fields, giving necessary and sufficient conditions for these polynomials to be irreducible over $\bF_q$.

In this paper, we study Fibonacci polynomials over finite fields that are perfect powers, analogous to the work of Cohn \cite{cohn1964square} and Bugeaud--Mignotte--Siksek \cite{bugeaud2006classical}.
In particular, we give a complete characterization of perfect powers.
\begin{theorem}
\label{thm:mainFn}
    Let $j > 1$. Let $p$ be an odd prime number coprime to \(2j\), $q$ be a power of $p$, and let $a = \mathrm{ord}_{2j}(p)$ denote the order of $p$ in $\left( \mathbb{Z}/2j\mathbb{Z} \right)^\times$. For $n > 0$, $F_n(T)$ is a perfect $j$-th power in $\mathbb{F}_q[T]$ if and only if $n = p^{ak}$ for some $k \in \mathbb{N}$.
\end{theorem}
Similar result is available for \(p = 2\) as well (Theorem \ref{thm:mainFnp2}).
Moreover, we give a characterization of Fibonacci polynomials that are \emph{powerful}, i.e. multiplicities of each irreducible factor are at least two.\footnote{Although we will mostly focus on nonconstant polynomials, constant polynomials will also be considered powerful, since they have no non-unit irreducible prime factors.}
\begin{theorem}
    \label{thm:mainFnpowerful}
    Let $p > 3$ be a prime and $\bF_q$ be a finite field of characteristic $p$.
    Then \(F_n(T)\) is powerful in \(\bF_q[T]\) if and only if \(p \mid n\).
\end{theorem}
Similar result holds for $p = 2$ and $p = 3$ as well (Corollary \ref{cor:powerfulp2} and \ref{cor:powerfulp3}).

These results contrast to the case of the original Fibonacci sequence.
There are only finitely many Fibonacci numbers that are perfect powers or powerful, where the second assertion is true under the ABC conjecture \cite{ribenboim1999abc,yabuta2007abc}.

In fact, we prove similar characterizations for certain \emph{generalized Lucas polynomial sequences} of Horadam \cite{Horadam1996}, which generalize both Fibonacci polynomials and Lucas polynomials (see Definition \ref{lucassq}).
The proofs of these characterizations are based on the polynomial analogue of Binet's formula and the discriminant formula of Florez, Higuita, and Ramirez \cite{florez2019resultant}, except that certain small characteristic cases need separate considerations.
We also provide explicit factorizations of these polynomials for certain $p$'s, where the accompanying code can be found in \href{https://github.com/seewoo5/sage-function-field}{github.com/seewoo5/sage-function-field}.

\subsection*{Acknowledgements}

This work was completed during the 2025 Berkeley Math REU program and supported by NSF RTG grant DMS-2342225, MPS Scholars, and the Gordon and Betty Moore Foundation.
We thank Tony Feng for organizing the REU.